% Lecture Template for ENGR 1120 030 031 E01- Tristan Hill - Spring 2017 - Spring 2018
% 
% Introduction to MATLAB 
%
% Intro to Logical and Branching

% Document settings
\documentclass[11pt]{article}
\usepackage[margin=1in]{geometry}
\usepackage[pdftex]{graphicx}
\usepackage{multirow}
\usepackage{setspace}
\usepackage{hyperref}
\usepackage{color,soul}
\usepackage{fancyvrb}
\usepackage{framed}
\usepackage{wasysym}
\usepackage{multicol}
\usepackage{array}

\pagestyle{plain}
\setlength\parindent{0pt}
\hypersetup{
    bookmarks=true,         % show bookmarks bar?
    unicode=false,          % non-Latin characters in Acrobat’s bookmarks
    pdftoolbar=true,        % show Acrobat’s toolbar?
    pdfmenubar=true,        % show Acrobat’s menu?
    pdffitwindow=false,     % window fit to page when opened
    pdfstartview={FitH},    % fits the width of the page to the window
    pdftitle={My title},    % title
    pdfauthor={Author},     % author
    pdfsubject={Subject},   % subject of the document
    pdfcreator={Creator},   % creator of the document
    pdfproducer={Producer}, % producer of the document
    pdfkeywords={keyword1} {key2} {key3}, % list of keywords
    pdfnewwindow=true,      % links in new window
    colorlinks=true,       % false: boxed links; true: colored links
    linkcolor=red,          % color of internal links (change box color with linkbordercolor)
    citecolor=green,        % color of links to bibliography
    filecolor=magenta,      % color of file links
    urlcolor=blue           % color of external links
}

% assignment number 
\newcommand{\NUM}{4 } 
\newcommand{\VSpaceSize}{2mm} 
\newcommand{\HSpaceSize}{2mm} 

\definecolor{mygray}{rgb}{.6, .6, .6}
\definecolor{mypurple}{rgb}{0.6,0.1961,0.8}
\definecolor{mybrown}{rgb}{0.5451,0.2706,0.0745}
\definecolor{mygreen}{rgb}{0, .39, 0}

\newcommand{\R}{\color{red}}
\newcommand{\B}{\color{blue}}
\newcommand{\BR}{\color{mybrown}}
\newcommand{\K}{\color{black}}
\newcommand{\G}{\color{mygreen}}
\newcommand{\PR}{\color{mypurple}}

\setulcolor{red} 
\setstcolor{green} 
\sethlcolor{mygray} 

\begin{document}

\textbf{ \LARGE ENGR 1120 Lecture Chapter \NUM - Selection Statements\\\\ Introduction to Logic} \\
	
\includegraphics[scale=0.3]{lecture1_fig1.png}\\\\

\Large
\begin{itemize}


	
	\item \textbf{ \LARGE What is \color{mypurple}Logic \color{black}? We need to learn some new vocab.}\\
	
		
		\begin{itemize}
			
			\item \textbf{\Large Types of Sentences (or expressions)}  
				\begin{itemize}
					\item Declarative \\
					\item Exclamatory\\
					\item Interogative\\
					\item Imperative\\
				\end{itemize}
			\item \textbf{ \Large What is a  \color{mypurple}Logical Expession \color{black}?} \\  \vspace{30mm}
            			
            		\item \textbf{ \Large a  \color{mypurple} Logical Expession \color{black}  is evaluated as ...} 
		\end{itemize}
\newpage

	\item \textbf{ \LARGE \color{mypurple}Logical Statments \color{black} are used for making {\it decisions}}\\
		\begin{itemize}
					\item This may be called {\it branching} or \color{blue} Selection \color{black} in programming \\
					\item Previously our programs were all {\it just a straight line}\\
					\item Solving advanced problems requires more complex \color{mybrown} Program Flow\color{black} \\

				\end{itemize}
				\newpage
				
					\item \textbf{ \LARGE \color{blue}Relational Operators \color{black} and \color{mypurple}Logical Expressions \color{black}}\\\\
		\begin{itemize}
					\item {\bf{\B Relational Operators}} combine {\bf{\G numerical expressions}} to make {\bf{\PR Logical Expressions }}.\vspace{20mm}\\
					\item The general form of a {\bf{\PR Logical Expression}} is shown below. \vspace{20mm}\\

						 \scalebox{1.1}{{\fontfamily{qcr}\selectfont Numerical  \hspace{7mm}Relational  \hspace{7mm}Numerical }} \\
						\scalebox{1.1}{{\fontfamily{qcr}\selectfont Expression \hspace{5mm} Operator \hspace{5mm} Expression }} \\\\
			
				\end{itemize}
				\newpage
				
					\item \textbf{ \LARGE \color{blue}if \color{black} Statments are used for making {\it decisions}.}\\
		\begin{itemize}
					\item \color{blue}if \color{black} is a reserved \color{mypurple} keyword \color{black}. You cannot use it for anything else. \\
					\item \color{blue}end \color{black} is a also a reserved \color{mypurple} keyword \color{black}. \\\\\\
					\item The general form of an \color{blue}If \color{black} statment is shown below. \vspace{20mm}\\ \\

						 \scalebox{1.1}{{\fontfamily{qcr}\selectfont  if (condition) }} \\\\
						 \scalebox{1.1}{{\fontfamily{qcr}\selectfont  \hspace{10mm}code }} \\\\
						 \scalebox{1.1}{{\fontfamily{qcr}\selectfont  end }} \\
				\end{itemize}
				\newpage
					\item \textbf{ {\BR Compound Logical Expressions} and {\G Logical Operators}} \\
			\begin{itemize}
				\item Logical Expressions can be combined to make Compound Logical Expressions\\\\
				
				\item This is done using {\bf{\G Logical Operators}}\\\\

				\item you are familiar with this idea...\\\\
				
				
					\LARGE
					\begin{center}
					\begin{tabular}{|| m{4cm} | m{5cm}| m{2.5cm} |m{3cm} || } \hline
						 {\bf English} &{\bf Comp. Science} & {\bf Symbol} & {\bf MATLAB} \\ \hline \hline
 						 &  &  & \\   
 						 &  &  & \\
 						 &  &  & \\
 						 &  &  & \\
 						 &  &  & \\\hline
 						 &  &  & \\
 						 &  &  & \\
 						 &  &  & \\
 						 &  &  & \\
 						 &  &  & \\\hline
 						 &  &  & \\
 						 &  &  & \\
 						 &  &  & \\
 						 &  &  & \\
 						 &  &  & \\ \hline
 						   
					\end{tabular}
					\end{center}


			
			\end{itemize}
				\newpage
	
				
					\item \textbf{We can represent this idea in a {\B Truth Table}}.\\
				
				\LARGE
					\begin{center}
					\begin{tabular}{|| m{2.75cm} | m{2.75cm}| m{2.75cm} |m{2.75cm}|m{2.75cm}|| } \hline
					 	 &  & &  &   \\
 						 &  & &  &   \\\hline\hline
 						 &  & &  &   \\
 						 &  & &  &   \\\hline
 						 &  & &  &   \\
 						 &  & &  &   \\\hline
 						 &  & &  &   \\
 						 &  & &  &   \\\hline
 						 &  & &  &   \\
 						 &  & &  &   \\\hline
 						 &  & &  &   \\
 						 &  & &  &   \\\hline
 						 &  & &  &   \\
 						 &  & &  &   \\\hline
 						 &  &  &  &   \\
 						 &  &  &  &   \\\hline
 						 &  &  &  &   \\
 						 &  &  &  &   \\\hline
 						   
					\end{tabular}
					\end{center}

				\newpage
	
				
					\item \textbf{We often represent the \color{mybrown} Program Flow \color{black} in a \color{mypurple} Flowchart\color{black} }\\
\end{itemize}


	

\end{document}



